\documentclass[10pt]{article}
\usepackage[margin=2.2cm]{geometry}
\usepackage{graphicx}
\usepackage{subcaption}
\usepackage{booktabs}
\usepackage{amsmath}
\usepackage{xcolor}
\usepackage[colorlinks=true, linkcolor=blue, citecolor=blue]{hyperref}
\graphicspath{{prl_comparison_figs/}}

\newcommand{\ok}{\textcolor{green!55!black}{\textbf{reproduced}}}
\newcommand{\partialok}{\textcolor{orange!85!black}{\textbf{partial}}}
\newcommand{\notok}{\textcolor{red!75!black}{\textbf{not reached}}}

\title{Side-by-side comparison: Duan, Agrawal \& Wang, PRL \textbf{134}, 048101 (2025)\\
vs.\ this codebase: (i) their SCFT reimplemented, (ii) charged CL-FTS}
\author{Internal research note (devel/charged\_polymers) --- figures from the cited paper\\
are reproduced here for private scholarly comparison only}
\date{2026-08-26}

\begin{document}
\maketitle

\begin{abstract}
Panel-by-panel comparison of the polyelectrolyte-brush results of
Duan, Agrawal \& Wang [PRL \textbf{134}, 048101 (2025)] --- obtained
with their electrostatic-correlation-augmented SCFT --- against our
independent reproduction with fully fluctuating charged complex-Langevin
FTS (smeared-charge compressible model, wall boxes, discrete chains).
Left panels: the paper's figures. Right panels: our measurements at the
strongest coupling inside our CL-stable envelope
($z_+{=}3$ ion correlation $u_3 \approx 13\,k_BT$, $\Theta$ solvent,
$\alpha = 1$ fully charged backbone, $N=25$, $\bar N = 10^4$).
Full provenance: \texttt{RESULTS\_FIG1\_REPRO.md} \S3--3f; all numbers
audited against the run data.
\end{abstract}

\section*{Part I --- their own theory, our implementation (correlation SCFT)}

Direct reimplementation of their variational/self-energy SCFT
(\texttt{prl\_scft.py}): repo discrete-chain propagator on a real-space
bond-convolution backend, guarded-Newton nonlinear PB with analytic
grand-canonical ions, LDA self-energies (SI Eq.~S33), exact $\xi$
elimination, $\alpha$- and correlation-strength continuation. All at
their exact Fig.~1 parameters.

\begin{figure}[h!]
\centering
\begin{subfigure}[t]{0.50\textwidth}
  \includegraphics[width=\linewidth]{prl_fig1}
  \caption{PRL Fig.~1.}
\end{subfigure}\hfill
\begin{subfigure}[t]{0.48\textwidth}
  \includegraphics[width=\linewidth]{ours_scft_fig1}
  \caption{Same theory, our implementation.}
\end{subfigure}
\end{figure}

\noindent\textbf{Verdict: \ok\ (quantitative).}
Curve-for-curve agreement, all states converged to $5\times10^{-4}$,
including the near-wall entropic depletion layer (hard-wall truncated
bond integral; graft source a Gaussian of width $0.5b$ centered on the wall
(half-Gaussian, the smooth analog of their $\delta(z)$; finite in the
discrete hard-wall formulation)): mean-field dashed family ($\phi\approx0.31$--$0.40$, extent
45--52\,nm, nearly valence-blind; theirs $0.32$--$0.38$, 45--50$b$);
neutral ($h=21.8$\,nm, plateau 0.51, edge $\sim$28$b$; theirs
$\sim$0.51, $\sim$28$b$); $z_+{=}2$ collapsed plateau $\phi=0.58$ to
18.7\,nm (theirs $\sim$0.59 to $\sim$19$b$) with $\Gamma=-0.056$ (their
algebraic Eq.~7 value: $-0.048$) and $\psi_S<0$; $z_+{=}3$ plateau
$\phi=0.72$--$0.73$, $h=14.5$\,nm (theirs $\sim$0.74 to $\sim$14$b$)
with $\Gamma=-0.32$ and $\psi_S=+0.66$ (charge inversion); $z_+{=}1$
slightly compacted, no overcharging. The Fig.~3 sign structure
($z_+{=}2$: overcharging without $\psi_S$ inversion; $z_+{=}3$: both)
is reproduced exactly. One physics note: the pointwise LDA self-energy
has no wavelength cutoff and goes grid-scale-unstable at $z_+{=}3$ full
coupling; smoothing the local ionic strength over the Born radius $a$
(the regularization their nonlocal $G_s$ carries intrinsically) restores
clean convergence and leaves $z_+{=}1,2$ untouched. Full account:
\texttt{RESULTS\_PRL\_SCFT.md}.

\subsection*{Their Figs.~2--4 with the reimplemented SCFT}

\begin{figure}[h!]
\centering
\begin{subfigure}[t]{0.49\textwidth}
  \includegraphics[width=\linewidth]{prl_fig2}
  \caption{PRL Fig.~2.}
\end{subfigure}\hfill
\begin{subfigure}[t]{0.49\textwidth}
  \includegraphics[width=\linewidth]{ours_scft_fig2}
  \caption{Ours, semicanonical with a GC monovalent background $c_1$
  (the physical form of the counterion reservoir; the onset position is
  reservoir-controlled --- shown by the $1/L$ box scaling --- and the
  paper reports neither box nor background). $c_1$ per valence aligns
  the ion-exchange crossings ($z_+{=}2$: $9.2\times10^{-4}$ vs their
  $\sim10^{-3}$\,M; $z_+{=}3$: $2.1\times10^{-5}$ vs their
  $\sim3\times10^{-5}$\,M); valleys $-0.42/-0.49$ (theirs
  $-0.45/-0.58$; the $z_+{=}3$ background screening shallows it). A
  single $c_1$ cannot align both (their crossing ratio 33 vs this
  tier's $\sim$14 --- the $I_{ex}$ signature).}
\end{subfigure}
\end{figure}

\begin{figure}[h!]
\centering
\begin{subfigure}[t]{0.49\textwidth}
  \includegraphics[width=\linewidth]{prl_fig3}
  \caption{PRL Fig.~3.}
\end{subfigure}\hfill
\begin{subfigure}[t]{0.49\textwidth}
  \includegraphics[width=\linewidth]{ours_scft_fig3}
  \caption{Ours. Thick blue: their Eq.~7 ($z_+{=}3$ min $-1.27$ vs
  their $-1.3$); thick orange: their Eq.~9 (their orange curve is this
  Donnan estimate, not a wall potential --- verified by shape). The
  $z_+{=}3$ peak: Eq.~9 with our $\Gamma$ gives $+3.9$, the true SCFT
  wall value $+0.9$; their $+2.5$ lies between and is NOT reproducible
  from their own published $\Gamma$ + Eq.~9 (which would give $+5$) ---
  an internal inconsistency traceable only to their unpublished
  numerics.}
\end{subfigure}
\end{figure}

\begin{figure}[h!]
\centering
\begin{subfigure}[t]{0.49\textwidth}
  \includegraphics[width=\linewidth]{prl_fig4}
  \caption{PRL Fig.~4.}
\end{subfigure}\hfill
\begin{subfigure}[t]{0.49\textwidth}
  \includegraphics[width=\linewidth]{ours_scft_fig4cd}
  \caption{Ours, their (c,d): $a_+ = 2.5, 2.0$\,\AA{} give smooth
  collapsed films; $a_+=1.5$\,\AA{} gives converged stationary
  oscillatory layers (period $\sim b$) --- the same physics as their
  panels, drawn as 1D profiles rather than density renders. Their
  (a,b): \textbf{the visual difference is dimensionality} --- their
  hexagonal pinned-micelle lattice is a 3D result, while our
  calculation (\texttt{prl\_scft2d.py}) is a 2D $(x,z)$ pilot in which
  micelles can only appear as stripe cross-sections by construction.
  The 2D pilot reproduces the lateral microphase separation itself:
  the uniform film is metastable and seeded runs converge to arrays of
  pinned dense domes ($\phi\sim0.75$, spacing 15--30\,nm, near-bare
  gaps; next figure). A full 3D extension would be required for the
  hexagonal arrangement.}
\end{subfigure}
\end{figure}

\begin{figure}[h!]
\centering
\includegraphics[width=0.9\textwidth]{ours_scft_fig4ab}
\caption*{Converged laterally structured state of the quenched 2D pilot
at their Fig.~4(a,b) conditions ($\sigma=0.03$\,nm$^{-2}$, $z_+{=}3$,
$\rho_b=1$\,mM): an array of pinned dense domes ($\phi\sim0.75$,
height $\sim$8\,nm, spacing 20\,nm, near-bare gaps), coexisting with
the metastable uniform film (first-order lateral microphase
separation).}
\end{figure}

\section*{Part II --- fully fluctuating charged CL-FTS vs the paper}

\textbf{Final (2026-08-29): three-way judgment at matched parameters,
hard-wall boxes on both sides.}
Their correlation SCFT (Part I implementation) applied to the CL
campaign's own system ($\Theta$, $N{=}25$, matched Born radii, matched
graft): mean fields agree to 1--6\%; the correlation DIFFERENTIALS show
the exact effect SATURATING
($\Delta h_1 = -0.095 \to -0.129 \to -0.131$ at $u_3 = 6.5/12/13\,k_BT$)
while the Gaussian LDA deepens toward its own fold
($-0.190 \to -0.291 \to$ no solution): overshoot $2.0\times$ at
$6.5\,k_BT$, $2.3\times$ at $12\,k_BT$, and beyond the runaway
threshold $u_3^* \approx 12.5\,k_BT$ the LDA electrolyte sector has no
physical solution (saturated-condensation branch) while CL remains
well-behaved. The LDA's overcharging ($\Gamma$ to $-0.25$) is
contradicted by three-seed-robust CL positives ($\Gamma \approx +0.09$)
at every tested coupling.

\begin{figure}[h!]
\centering
\includegraphics[width=0.85\textwidth]{cl_vs_scft_fig1}
\caption*{LDA (red, deterministic scan $u_3 = 6.5$--$12\,k_BT$) vs
exact CL-FTS (blue, 3 seeds), same hard-wall microscopic system.
Dotted line: LDA runaway threshold ($u_3^* \approx 12.5\,k_BT$),
beyond which their theory has no physical solution.}
\end{figure}


\section*{Fig.~1 --- density profiles: valence-selective correlation collapse}

\begin{figure}[h!]
\centering
\begin{subfigure}[t]{0.50\textwidth}
  \includegraphics[width=\linewidth]{prl_fig1}
  \caption{PRL Fig.~1: with correlations (solid), the $z_+{=}3$ brush
  collapses \emph{below the neutral profile}; mean field (dashed) is
  swollen and nearly $z_+$-insensitive. ($N{=}100$, $b{=}1$\,nm,
  $a{=}2.5$\,\AA, $I_b{=}0.3$\,M.)}
\end{subfigure}\hfill
\begin{subfigure}[t]{0.46\textwidth}
  \includegraphics[width=\linewidth]{ours_fig1}
  \caption{Ours, HARD-WALL rerun ($\Theta$, $u_3{=}12\,k_BT$, 3 seeds;
  dashed: quasi-saddle mean field, solid: CL): profiles now show the
  correct wall depletion layer ($\phi(0){\approx}2{\times}10^{-4}$).
  Correlations collapse the $z_+{=}3$ brush ($h_1$: $0.80 \to 0.67$)
  while $z_+{=}1$ barely moves; the collapsed brush remains
  \emph{above} our neutral reference ($h_1{=}0.62$).}
\end{subfigure}
\end{figure}

\noindent\textbf{Verdict: \ok\ (mechanism).} The defining structure ---
correlation-induced, $z_+$-selective collapse invisible at mean field ---
appears identically. The causality was pinned three ways
(\S3, reflecting-wall era; mechanism unchanged by the wall fix): the
contraction scales as $z_+^2$ (ratio $8.7 \approx 9$), vanishes in the
neutral control, and tracks the Coulomb coupling $E$.
\textbf{Difference, resolved in \S3h:} the paper's collapse undershoots
the neutral height; ours saturates well above it. The matched-parameter
comparison shows this is not a limitation of our envelope --- the exact
effect saturates while their Gaussian LDA keeps deepening ---
so the below-neutral amplitude is an artifact of their approximation.

\section*{Fig.~2 --- nonmonotonic $h(\rho_b)$: collapse and reexpansion}

\begin{figure}[h!]
\centering
\begin{subfigure}[t]{0.50\textwidth}
  \includegraphics[width=\linewidth]{prl_fig2}
  \caption{PRL Fig.~2: $\log(h/h_0)$ vs $\rho_b$. For $z_+{=}2,3$ the
  height is nonmonotonic --- collapse below the neutral line, then
  reexpansion --- matching the Tirrell-group experiments (symbols).}
\end{subfigure}\hfill
\begin{subfigure}[t]{0.46\textwidth}
  \includegraphics[width=\linewidth]{ours_fig2}
  \caption{Ours, on THEIR axis and axes (monovalent counterions + 1:1
  background salt, multivalent SALT scanned; hard wall, 3D; x = the
  MEASURED reservoir concentration of the canonical box,
  $b \equiv 1$\,nm): exact CL valence series --- common extended
  plateau, valence-ordered collapse $z_+{=}3 > 2 > 1$ ($z_+{=}1$
  tracks mean field), floor above the neutral brush (dashed),
  nonmonotonic $\Gamma$ (\S3i-d,g).}
\end{subfigure}
\end{figure}

\begin{figure}[h!]
\centering
\includegraphics[width=0.8\textwidth]{ours_fig2_threeway}
\caption*{Fig.~2-1 (FINAL) --- the $z_+{=}3$ three-way judgment at
identical parameters ($N{=}50$ brush, $\sigma{=}30$, tall box,
$u_3 = 12.9\,k_BT$; 8 salt points, 2 seeds). Mean field (gray):
monotone. Their PRL SCFT (orange): dives to $-0.44$ at
${\sim}10^{-2}$\,M --- essentially its neutral line --- with
overcharging and reexpansion. Exact FTS (green): smooth collapse to
$-0.30$, sitting BETWEEN mean field and their theory at every
measurable point (the exact correlation deepening is $40$--$50$\% of
their prediction), staying $0.17$ above the exact neutral brush
(dotted, $-0.47$) with $\Gamma > 0$ throughout ($+0.67 \to +0.11$).
Left-arrows: onset-region points whose bulk trivalent is fully
depleted by the brush (canonical box; reservoir concentration below
detection). The $N{=}25$ version is in \S3i-g.}
\end{figure}

\begin{figure}[h!]
\centering
\includegraphics[width=0.78\textwidth]{ours_fig2_cfamily}
\caption*{Fig.~2-2 --- interpolating their correction, $F = F_{\rm MF}
+ c\,F_{\rm 1loop}$: the exact ensemble (green) lies on the
$c{=}0.5$ member of the family across the whole collapse region.
At $u_3 = 12.9\,k_BT$ the non-Gaussian terms cancel HALF of the
Gaussian correlation free energy: no physical parameter maps their
theory onto exact sampling --- only the truncation weight $c \approx
0.5$ does (\S3i-i).}
\end{figure}

\begin{figure}[h!]
\centering
\includegraphics[width=0.62\textwidth]{cstar_vs_u3}
\caption*{Fig.~2-3 --- two-observable $c^*(u_3)$: the fraction of
their Gaussian correction confirmed by exact sampling, extracted
independently from brush height ($c^*_h$, purple) and condensation
($c^*_\Gamma$, brown) at five coupling rungs (fixed $s{=}0.02$).
Nowhere is the full correction ($c^*{=}1$) supported. At weak coupling
both observables give the same $c^* \approx 0.6$ (a genuine effective
theory); beyond $u_3 \sim 6\,k_BT$ they SPLIT ($0.46$ vs $0.24$ at
$u_3{=}12$--$13$): no single rescaling fixes both observables --- the
Gaussian truncation carries a form error, consistent with
non-Gaussian ion-bridging attraction absent from any single-ion
self-energy (\S3i-j).}
\end{figure}

\noindent\textbf{Verdict: \ok\ (mechanism, on the correct axis) ---
amplitude gap attributed to the LDA (\S3i-b--3i-d).} An important
correction: our earlier salt scans used trivalent brush COUNTERIONS, an
axis on which the brush starts collapsed and no valley can exist by
construction. On the paper's actual axis (monovalent counterions,
trivalent ions arriving only via salt) the exact ensemble DOES
reproduce the Fig.~2 phenomenology: extended $\to$ collapse (
$h_1$: $1.07 \to 0.70$, CL below MF throughout, gap growing with salt)
$\to$ valley floor, with nonmonotonic $\Gamma$. What does NOT appear
exactly: the amplitude. Their collapse undershoots the neutral brush by
${\sim}40$\%; our exact valley floor stays above ours. Attribution
(\S3i-g, magenta curve above): their Gaussian LDA, run on our system
at the matched measured ionic environment, gives the full deep valley
--- onset one decade earlier, $1.75\times$ deeper, below neutral, with
an overcharging window --- none of which survives exact sampling. (On
the trivalent-counterion axis, figure below, the theory additionally
loses its solution entirely at low screening.)

\begin{figure}[h!]
\centering
\includegraphics[width=0.72\textwidth]{lda_overlay}
\caption*{Attribution overlay (trivalent-counterion axis, hard wall,
$u_3{=}12\,k_BT$): their LDA (red) dives below its neutral line and
loses its solution (runaway $\times$); exact CL (blue) saturates at its
neutral floor.}
\end{figure}

\section*{Fig.~3 --- effective brush charge $\Gamma$: neutralization vs overcharging}

\begin{figure}[h!]
\centering
\begin{subfigure}[t]{0.50\textwidth}
  \includegraphics[width=\linewidth]{prl_fig3}
  \caption{PRL Fig.~3: for $z_+{=}2,3$, $\Gamma$ dips \emph{below zero}
  in a salt window (overcharging), before recovering at high salt.}
\end{subfigure}\hfill
\begin{subfigure}[t]{0.46\textwidth}
  \includegraphics[width=\linewidth]{ours_fig3}
  \caption{Ours, HARD-WALL rerun: correlations drive $\Gamma$ from the
  mean-field $0.24$--$0.30$ (dashed) to a $0^+$ plateau (complete
  neutralization, $+0.006$--$0.010$ at 150k steps) across the low-salt window,
  rising to $+0.14$ at high salt. No blue point crosses zero. Halving
  the ion smearing (orange, hard-wall, 150k, 2 seeds) pushes the
  low-salt point to $-0.017 \pm 0.017$ --- consistent with complete
  neutralization, at most marginal overcharging at doubled per-ion
  coupling.}
\end{subfigure}
\end{figure}

\noindent\textbf{Verdict: \partialok, revised by the hard-wall reruns
(\S3h--3i).} Correlation-driven condensation far beyond mean field is
real and now hard-wall-confirmed: $\Gamma$ drops from the mean-field
$0.24$--$0.30$ to a $0^+$ plateau ($+0.008$--$0.012$) in the low-salt
window (\S3i), and at fixed $s{=}0.02$ the exact ensemble SATURATES at
$\Gamma \approx +0.09$ across $u_3=6.5$--$13\,k_BT$ (\S3h). It never
crosses zero. The paper's overcharging window ($\Gamma<0$) is
where their Gaussian LDA is already overshooting ours by
$2.0$--$2.3\times$ and heading into its runaway ($u_3^* \approx
12.5\,k_BT$); we now attribute their $\Gamma<0$ amplitude to the
Gaussian/LDA approximation rather than to physics our envelope
``cannot yet reach.''

\section*{Fig.~4 --- lateral microphase separation (precursor only)}

\begin{figure}[h!]
\centering
\begin{subfigure}[t]{0.46\textwidth}
  \includegraphics[width=\linewidth]{prl_fig4}
  \caption{PRL Fig.~4: pinned micelles in hexagonal arrangement at
  moderate grafting density; oscillatory layers at high density.}
\end{subfigure}\hfill
\begin{subfigure}[t]{0.50\textwidth}
  \includegraphics[width=\linewidth]{ours_fig4}
  \caption{Ours, HARD-WALL rerun ($s{=}0.001$, three seeds, two times
  each): instantaneous $P$ column-density snapshots show box-scale
  lateral clustering (relative std $0.36$--$0.44$, single cluster per
  box) --- the precursor survives the wall fix --- but our lateral box
  ($0.8\,R_0$) is smaller than the expected micelle spacing, so no
  lattice can form. A $3$--$4\,R_0$ lateral box is the designed
  follow-up.}
\end{subfigure}
\end{figure}

\noindent\textbf{Verdict: precursor observed, geometry-limited.}
The lateral-inhomogeneity indicator rises to $\sim$10$\times$ the neutral
baseline for charged brushes, strongest at low salt (their Fig.~4 lives
at $\rho_b \sim 1$\,mM), but a quantitative Fig.~4 comparison requires
the wide-box campaign.

\section*{Correspondence of models and scales}

\begin{table}[h!]
\centering\small
\begin{tabular}{lll}
\toprule
 & PRL 134, 048101 & this work \\
\midrule
theory & variational Gaussian (self-energy) SCFT & full CL-FTS sampling \\
fluctuations & one-loop, self-consistent & all orders (stochastic) \\
chain & discrete, $N=100$, $\alpha=1$ & discrete, $N=25$, $\alpha=1$ \\
solvent & $\Theta$ ($\chi = 0.5$) & $\Theta$ ($\chi N = 12.5$) \\
ions & explicit, Born radius 2.5\,\AA, grand-canonical & explicit 1-segment, smeared
 ($a_{\rm ion}=0.025$--$0.0125\,R_0$), canonical \\
$z_+{=}3$ ion correlation & $\sim 7\,k_BT$ (near point-like) & $\leq 13\,k_BT$ nominal,
 cutoff-limited \\
geometry & 1D--3D wall, analytic-ion Poisson & 3D all-reflecting box, fluctuating $\psi$ (DCT/ETD) \\
\bottomrule
\end{tabular}
\end{table}

\section*{Summary scoreboard}

\begin{table}[h!]
\centering\small
\begin{tabular}{lll}
\toprule
PRL result & status & evidence \\
\midrule
\multicolumn{3}{l}{\emph{Part I: their SCFT, our implementation}} \\
Fig.~1 all profiles (MF, corr, neutral) & \ok\ (quantitative) & curve-for-curve match \\
Fig.~1--2 collapse below neutral & \ok & $z_+{=}2$: 18.7, $z_+{=}3$: 14.5 vs neutral 21.8\,nm \\
Fig.~3 overcharging $\Gamma<0$ ($z_+{=}2,3$) & \ok & $-0.057$ / $-0.32$; Eq.~7: $-0.048$ / $-0.65$ \\
Fig.~3 $\psi_S$ inversion only at $z_+{=}3$ & \ok & $\psi_S$: $-0.81$ / $-0.15$ / $+0.75$ \\
\midrule
\multicolumn{3}{l}{\emph{Part II: fully fluctuating CL-FTS (honest restatement, 2026-08-29;}} \\
\multicolumn{3}{l}{\emph{\S3--3f were REFLECTING-wall runs, superseded by the hard-wall reruns of \S3h--3i)}} \\
Fig.~1 contraction mechanism & \ok\ (mechanism only) & $z_+^2$ ratio, $E$-scaling, neutral null; small amplitude (\S3, 3c, 3h) \\
Fig.~2 collapse mechanism (correct axis) & \ok & valence series $z_+{=}1/2/3$: common plateau, ordered collapse, valley floor + nonmonotonic $\Gamma$; CL $<$ MF throughout (\S3i-d,f) \\
Fig.~2 below-neutral valley depth & \notok\ (attributed) & matched ionic env.: their theory collapses 1 decade earlier, $1.75\times$ deeper, below neutral, $\Gamma<0$; exact does none (\S3i-g) \\
Fig.~3 strong condensation & \ok & hard-wall: $\Gamma$: MF $0.24$--$0.30 \to 0^+$ plateau at low salt; saturates $+0.09$ vs coupling (\S3h--3i) \\
Fig.~3 overcharging ($\Gamma<0$) & \textcolor{red!75!black}{\textbf{refuted}} & exact ensemble stays positive at all tested couplings; LDA's $\Gamma<0$ is a Gaussian artifact (\S3h) \\
Fig.~1--2 collapse below neutral & \textcolor{red!75!black}{\textbf{refuted at matched params}} & exact contraction saturates far above the paper's amplitudes (\S3h) \\
Fig.~4 lateral micelles & \partialok & box-scale precursor survives hard wall (\S3i); dome arrays in 2D SCFT only \\
Their theory vs exact sampling & --- & LDA overshoots $2.0$--$2.3\times$, runaway at $u_3^* \approx 12.5\,k_BT$ (\S3h) \\
\bottomrule
\end{tabular}
\end{table}

\noindent The amplitude gap traces to a single, well-characterized cause:
our UV-regularized (smeared-charge) model caps the ion-correlation
strength that their near-point-charge treatment attains, and shrinking the
smearing collides with a CL stability boundary. Candidate paths beyond it
--- their grand-canonical analytic-ion formulation (SI Eq.~S24) grafted
onto our fluctuating-$\psi$ sector, or a dedicated strong-coupling
$\psi$ stabilization --- are recorded as future work.

\end{document}
